\section{Lattice Yang-Mills Theory: Formulation}
\label{sec:lattice}

We define the theory on a hypercubic lattice $\Lambda = (a\mathbb{Z})^4$ with lattice spacing $a$.

\subsection{Configuration Space}
The gauge field is represented by link variables $U_{x,\mu} \in SU(N)$. The configuration space is the product of group manifolds $\prod_{l \in \Lambda} SU(N)$.

\subsection{Action and Measure}
The dynamics are governed by the Wilson action:
\begin{equation}
S_W(U) = \beta \sum_{p} \left( 1 - \frac{1}{N} \text{Re Tr}(U_p) \right)
\end{equation}
where $U_p$ is the path-ordered product of link variables around a plaquette $p$, and $\beta = 2N/g^2$ is the inverse coupling.
The expectation value of an observable $\mathcal{O}$ is given by the Haar measure integral:
\begin{equation}
\langle \mathcal{O} \rangle = \frac{1}{Z} \int \prod_{l} dU_l \, \mathcal{O}(U) \, e^{-S_W(U)}
\end{equation}

\subsection{Transfer Matrix and Hamiltonian}
The existence of a positive Hamiltonian $H \ge 0$ is guaranteed by the property of Reflection Positivity. The transfer matrix $T = e^{-aH}$ describes evolution in the Euclidean time direction. A strictly positive mass gap $\Delta > 0$ corresponds to the condition that the spectrum of $H$ (excluding the vacuum 0) is bounded below by $\Delta$.
In any finite volume $V$, with compact gauge group $SU(N)$, the spectrum of $H$ is discrete and the gap $\Delta(V) > 0$ is trivially non-zero. The physical challenge is to prove $\lim_{V \to \infty} \Delta(V) > 0$.
